\begin{frame}{Detecção de Comunidades}
    \framesubtitle{Algoritmo de Louvain}
    \begin{block}{Objetivo}
        Identificar agrupamentos estruturais de atletas com padrões competitivos similares.
    \end{block}

    \begin{columns}[T]
        \begin{column}{0.48\textwidth}
            \textbf{Resultados Gerais:}
            \begin{itemize}
                \item \textbf{36 comunidades} detectadas
                \item \textbf{4.659 atletas} analisados
                \item \textbf{0 nós isolados} em todas as redes
                \item Modularidade varia por esporte
            \end{itemize}
        \end{column}
        \begin{column}{0.48\textwidth}
            \textbf{Distribuição:}
            \begin{itemize}
                \item Swimming: 26 comunidades
                \item Basketball: 5 comunidades
                \item Football: 5 comunidades
            \end{itemize}
        \end{column}
    \end{columns}
\end{frame}

\begin{frame}{Modularidade por Modalidade}
    \framesubtitle{Diferenças Estruturais entre Esportes}
    \begin{table}[h]
        \centering
        \begin{tabular}{lccc}
            \toprule
            \textbf{Esporte} & \textbf{Tipo} & \textbf{Modularidade} & \textbf{Interpretação} \\
            \midrule
            Swimming & Individual & 0.73 & Alta segregação \\
            Swimming & Team & 0.45 & Média segregação \\
            Basketball & Coletivo & 0.003 & Estrutura coesa \\
            Football & Coletivo & 0.0006 & Estrutura coesa \\
            \bottomrule
        \end{tabular}
    \end{table}

    \vspace{0.3cm}
    \begin{block}{Interpretação}
        \begin{itemize}
            \item \textbf{Alta modularidade} (Swimming): Comunidades bem definidas por especialização técnica, período temporal e concentração geográfica
            \item \textbf{Baixa modularidade} (Basketball/Football): Rede densamente interconectada com fronteiras difusas entre comunidades
        \end{itemize}
    \end{block}
\end{frame}

\begin{frame}{Natureza Multifatorial das Comunidades}
    \framesubtitle{Além de Geografia e Tempo}
    \begin{alertblock}{Achado Principal}
        Comunidades NÃO correspondem a agrupamentos simples por período temporal ou nacionalidade, mas emergem de \textbf{combinações complexas} de múltiplos fatores.
    \end{alertblock}

    \vspace{0.3cm}
    \begin{columns}[T]
        \begin{column}{0.48\textwidth}
            \textbf{Natação:}
            \begin{itemize}
                \item Diversidade temporal moderada
                \item Entropia média: 1.60
                \item Concentração geográfica moderada
                \item Comunidades especializadas
            \end{itemize}
        \end{column}
        \begin{column}{0.48\textwidth}
            \textbf{Esportes Coletivos:}
            \begin{itemize}
                \item Alta diversidade temporal
                \item Entropia: 2.46-2.53
                \item Distribuição geográfica equilibrada
                \item Comunidades multi-era
            \end{itemize}
        \end{column}
    \end{columns}

    \vspace{0.2cm}
    \begin{block}{Conclusão}
        \footnotesize
        Algoritmos de detecção capturam estruturas emergentes não redutíveis a características isoladas.
    \end{block}
\end{frame}

\begin{frame}{Perfis de Medalhas por Comunidade}
    \framesubtitle{Distribuição Aproximadamente Equilibrada}
    \begin{block}{Achado Principal}
        Distribuição de medalhas nas 36 comunidades apresenta padrão predominantemente equilibrado entre os três tipos.
    \end{block}

    \begin{table}[h]
        \centering
        \footnotesize
        \begin{tabular}{lcccc}
            \toprule
            \textbf{Estatística} & \textbf{Ouro (\%)} & \textbf{Prata (\%)} & \textbf{Bronze (\%)} & \textbf{N} \\
            \midrule
            Média         & 31.8      & 33.4       & 34.8        & 36 \\
            Mínimo        & 16.7      & 20.0       & 25.0        & -- \\
            Máximo        & 50.0      & 50.0       & 50.0        & -- \\
            \bottomrule
        \end{tabular}
    \end{table}

    \vspace{0.3cm}
    \begin{block}{Interpretação}
        \footnotesize
        Distribuição aproximadamente uniforme (31.8\%/33.4\%/34.8\%) caracteriza competição sistemática sem monopolização absoluta de posições superiores no pódio.
    \end{block}
\end{frame}
